\subsection{版本、勘误与复算入口}
\phantomsection\label{app:revision-notes}

\noindent\textbf{版本：}\bookbuildsource{}复算时先核对输入、形状和计算规则，再确认模型版本与依赖。\par
\noindent\textbf{勘误：}反馈时请注明章节、输入与期待结果；已确认的问题见仓库\href{https://github.com/CoCoThink/ai-by-intuition/blob/main/CONTRIBUTING.md\#\%E5\%8B\%98\%E8\%AF\%AF\%E4\%B8\%8E\%E5\%8F\%8D\%E9\%A6\%88}{《编写与维护指南》的“勘误与反馈”小节}，反馈入口为 Issues。\par
\noindent\textbf{复现：}全部脚本、依赖与预期输出见\hyperref[app:reproduction-scripts]{附录~\ref*{app:reproduction-scripts} 的复算脚本表（第~\pageref*{app:reproduction-scripts}~页）}；复现本 PDF 时切换到上面的提交。\par
\noindent\url{https://github.com/CoCoThink/ai-by-intuition}\par

% BEGIN GENERATED DATED CLAIMS
\subsubsection*{带日期断言清单}
\phantomsection\label{app:dated-claims}
下表列出带日期的断言：所用版本、内容核验日（UTC）与核验范围；来源均列于附录 F。
核验日指实际内容复核日期；重新编译不会刷新核验日期。
\begingroup\small\setlength{\tabcolsep}{3pt}\renewcommand{\arraystretch}{1.25}
\tagpdfsetup{table/header-rows={1}}
\begin{longtable}{@{}>{\raggedright\arraybackslash}p{50mm}>{\raggedright\arraybackslash}p{26mm}>{\raggedright\arraybackslash}p{80mm}@{}}
\toprule 断言与正文位置 & 版本／核验日（UTC） & 来源与核验范围\\\midrule
\endfirsthead
\toprule 断言与正文位置 & 版本／核验日（UTC） & 来源与核验范围\\\midrule
\endhead
\bottomrule\endfoot
MCP 移除初始化与协议级会话；每次请求携带版本和客户端能力（\hyperref[concept:mcp]{第~\pageref*{concept:mcp}~页}） & 2026-07-28\newline 2026-10-25 & \href{https://modelcontextprotocol.io/specification/2026-07-28/basic/index}{核心消息}、\href{https://modelcontextprotocol.io/specification/2026-07-28/changelog}{变更记录}、\href{https://modelcontextprotocol.io/specification/2026-07-28/server/tools}{工具调用}。已阅读规范并核对正文请求体；HTTP 传输头另见规范。\\
Sampling、Roots、Logging 与 OAuth 动态客户端注册已弃用，兼容窗口内仍可用（\hyperref[concept:mcp]{第~\pageref*{concept:mcp}~页}） & 2026-07-28\newline 2026-11-06 & \href{https://modelcontextprotocol.io/specification/2026-07-28/changelog}{官方变更记录的 Deprecated 清单}。已阅读弃用清单与迁移说明；区分弃用窗口内可用的特性与已移除的初始化和协议级会话。\\
MCP 宿主管理客户端；每个客户端与一个服务端通信（\hyperref[concept:mcp]{第~\pageref*{concept:mcp}~页}） & 2026-07-28\newline 2026-10-28 & \href{https://modelcontextprotocol.io/specification/2026-07-28/architecture/index}{架构规范}。已核对宿主授权与编排、客户端一对一通信和服务端能力；模型请求由应用转成协议消息。\\
Claude 工具定义含名称、描述、输入模式；按停止原因继续工具循环（\hyperref[elective:real-function-call]{第~\pageref*{elective:real-function-call}~页}） & 2026-10-21\newline 2026-10-25 & \href{https://platform.claude.com/docs/en/api/messages/create}{Messages}、\href{https://platform.claude.com/docs/en/api/overview}{认证与请求头}、\href{https://platform.claude.com/docs/en/agents-and-tools/tool-use/define-tools}{定义工具}、\href{https://platform.claude.com/docs/en/agents-and-tools/tool-use/handle-tool-calls}{处理调用}。核验方式为官方规范对照与离线多轮闭环，覆盖必填请求体、认证与版本头。\\
A2A 的消息、任务与产物交换机制（\hyperref[concept:a2a]{第~\pageref*{concept:a2a}~页}） & v1.0.0\newline 2026-10-25 & \href{https://a2a-protocol.org/v1.0.0/specification/}{固定版本规范}。核验对象为本书采用的 v1.0.0 规范；核验方式为固定版本的字段与流程对照。\\
GPT-2 small：12 层、12 头、隐藏宽度 768，约 1.244 亿参数；GELU、可学习 LayerNorm 与共享输出头（\hyperref[intro:23a-from-toy-to-gpt]{第~\pageref*{intro:23a-from-toy-to-gpt}~页}） & 607a30d\newline 2026-10-29 & \href{https://huggingface.co/openai-community/gpt2/blob/607a30d/config.json}{固定配置}、\href{https://github.com/openai/gpt-2/blob/c2dae27c1029770cea40978813f17a5fd545b883/src/model.py}{固定模型实现}。本日重新核对固定配置与 c2dae27 模型实现：词表、位置范围、FFN 宽度、GELU、归一化、偏置及输出头共享；按可训练元素复算 124,439,808 参数。权重推理仍引用 2026-10-23 存档实测。\\
Qwen2.5-0.5B：24 层、14 个 Query 头、2 个 KV 头、头宽 64（\hyperref[example:gqa-budget]{第~\pageref*{example:gqa-budget}~页}） & 060db64\newline 2026-10-30 & \href{https://huggingface.co/Qwen/Qwen2.5-0.5B/blob/060db64/config.json}{固定配置}。已读取四个配置字段，复算 R-13 与 32K 缓存；预算不含运行时开销。\\
真实 Qwen 分词器对“春天花开了。”的切分与 ID（\hyperref[experiment:real-chinese-tokenizer]{第~\pageref*{experiment:real-chinese-tokenizer}~页}） & 060db64\newline 2026-10-25 & \href{https://huggingface.co/Qwen/Qwen2.5-0.5B/blob/060db6499f32faf8b98477b0a26969ef7d8b9987/tokenizer.json}{固定分词器}。已运行固定分词器；只复算分词，不加载该模型权重。\\
DeepSeek-V4.1-Flash 的因果编码器—解码器：40 层分为 20+20，输入／输出每 Token 激活 8B／16B（\hyperref[elective:causal-encoder-decoder]{第~\pageref*{elective:causal-encoder-decoder}~页}） & 2026-09-10 公告／模型卡\newline 2026-10-25 & \href{https://huggingface.co/deepseek-ai/DeepSeek-V4.1-Flash}{DeepSeek 模型卡}、\href{https://www.deepseek.com/en/news/deepseek-v4-1-flash/}{DeepSeek 公告}。已核对厂商公开的结构描述；速度与基准成绩的来源为厂商公布的测评记录。\\
LLM 评审存在位置、冗长与自我偏好等偏差（\hyperref[concept:llm-judge]{第~\pageref*{concept:llm-judge}~页}） & 2023 论文\newline 2026-10-25 & \href{https://arxiv.org/abs/2306.05685}{原论文}。已核对论文所列偏差；正文要求交换答案顺序并以人工样本校验。\\
PyTorch nn.Dropout：训练时独立置零并缩放，评估时直接传递（\hyperref[concept:dropout]{第~\pageref*{concept:dropout}~页}） & 2.14 文档\newline 2026-10-27 & \href{https://docs.pytorch.org/docs/2.14/generated/torch.nn.Dropout.html}{Dropout}、\href{https://docs.pytorch.org/docs/2.14/generated/torch.nn.Module.html}{评估模式}。核验方式为该版官方文档对照与缩放例子的手算复核。\\
\end{longtable}\endgroup
% END GENERATED DATED CLAIMS

\subsubsection*{真实中文分词：固定修订与实测记录}
\phantomsection\label{app:tokenizer-record}

第~\ref*{intro:20a-tokenizer}~章的中文切分使用 Qwen2.5-0.5B 分词器，
修订为 \nolinkurl{060db6499f32faf8b98477b0a26969ef7d8b9987}，
设置 \texttt{add\_special\_tokens=False}。
分词器文件哈希和实测记录见仓库 \path{scripts/fixtures/qwen25_tokenizer_trace.json}。

\subsubsection*{真实 GPT-2：运行方式与实测来源}

\texttt{python3 scripts/gpt2\_demo.py}
不加载权重时复核概率公式。安装兼容的 \texttt{torch} 和 \texttt{transformers} 后，
用 \texttt{-{}-real} 加载固定版本的 GPT-2 small，核对 Token、logits 与概率。

\texttt{-{}-local-files-only} 仅读本地缓存，安装 Matplotlib 后可用
\texttt{-{}-heatmap} 保存注意力图。复现时保留输入、权重修订和环境信息。

正文的 GPT-2 small 实测由输入、权重修订和运行环境共同确定；
完整依赖、运行记录与热图见仓库\href{https://github.com/CoCoThink/ai-by-intuition/blob/main/CONTRIBUTING.md\#gpt-2-\%E7\%9C\%9F\%E5\%AE\%9E\%E6\%9D\%83\%E9\%87\%8D\%E5\%AE\%9E\%E9\%AA\%8C}{《编写与维护指南》的“GPT-2 真实权重实验”小节}。

\textbf{2026-10-23 已存档实测（GPT-2 small，$T=1.0$）：}
Prompt：\texttt{The capital of France is}。
输入 Token ID：\texttt{[464, 3139, 286, 4881, 318]}。
logits 形状为 $1\times5\times50{,}257$，概率和为 $1.00000000$。
Top-5（Token ID ／ 解码文本 ／ 概率）：

\begin{center}\small
\tagpdfsetup{table/header-rows={1}}
\begin{tabularx}{\linewidth}{@{}r>{\raggedright\arraybackslash}Xr@{}}
\toprule
ID & 解码文本 & 概率\\\midrule
262 & \texttt{' the'} & 0.08459\\
783 & \texttt{' now'} & 0.04795\\
257 & \texttt{' a'} & 0.04616\\
4881 & \texttt{' France'} & 0.03238\\
6342 & \texttt{' Paris'} & 0.03224\\
\bottomrule
\end{tabularx}
\end{center}
表中列出这个前缀之后各候选 Token 的概率；综合能力的评测方法见第~\ref*{intro:23e-tech-judgement}~章。
